\documentclass[FGBook.tex]{subfiles}

\begin{document}

\proza{Proměna}{FRANZ KAFKA}

\noindent\begin{wrapfigure}{r}{0.25\textwidth}\footnotesize
\subwection{Ukázka z díla:}
\setlength{\parindent}{3pt}
Největší okamžité zlepšení situace jim ovšem jistě snadno přinese změna bytu;
vezmou si teď menší a levnější a přitom lépe položený a vůbec praktičtější byt,
než je nynější, který vybíral ještě Řehoř.
Jak si tak povídali, zadívali se pan a paní Samsovi na svou čím dál čilejší dceru
a skoro zároveň je napadlo, jak za poslední dobu přes všechno soužení,
od něhož jí pobledly tváře, rozkvetla v~krásnou a kyprou dívku.
Umlkli, a dorozumívajíce se skoro nevědomky pohledy, pomysleli si,
že bude teď načase, aby pro ni také vyhledali hodného muže.
A~připadalo jim jako dotvrzení jejich nových snů a dobrých úmyslů,
když u cíle jejich cesty dcera prví vstala a protáhla své mladé tělo.
\end{wrapfigure}

\wection{LITERÁRNÍ TEORIE}

\subwection{Literární druh a žánr:}
\noindent epická próza, \texthl{novela}

\subwection{Literární směr:}
\noindent Kafkovo dílo je výjimečné.
Pochází z prostředí, které prolínalo vícero směrů a Kafka to dotáhnul o úroveň výše.
I když se dílo dá řadit do \textbfhl{expresionismu}, tak k jeho motivům více sedí \textbfhl{existencialismus}.
Jedná~se~o~dílo, jenž vyjadřuje vnitřní pocity muže,
který se i díky \textithl{odcizení} od sebe samotného ocitl v~této~absurdní situaci.
A díky své proměně se odcizil i od své rodiny a nezbylo mu nic jiného, než uvažovat sám nad sebou.

\subwection{Jazyk a slovní zásoba:}
\noindent Spisovný a formální jazyk.
Věcný k popisu tak absurdních situací.
Autor svojí slovní zásobou popisuje \textithl{pocity odcizení, úzkosti a bezmoci}.

% \subwection{Figury:}

% \subwection{Tropy:}

\subwection{Stylistická charakteristika textu:}
\noindent Strohý a detailní popis hmyzího těla a Řehořova pokoje.
\texthl{Absurdní situace} věrná Kafkově jménu, popsána věcně, jako by se nic nedělo.

\subwection{Postavy:}
\noindent 
\textschl{Řehoř Samsa --} pracovitý mladý muž,
jen pracuje aby \textithl{uživil svojí rodinu}.
Absurdně se z něj \textithl{stane škůdce},
rodina na něj postupně zanevře a~Řehoř umírá. \\
\textschl{Matka --} slabá ale milující matka.
Řehoře se dlouhou dobu zastává, dokud se nepřesvědčí,
že to stvoření \texthl{není její syn}. \\
\textschl{Otec --} na Řehoře zapomíná jako první,
kvůli synově neschopnosti se musí vrátit do práce.
I když Řehoře prvně do práce nutil~s~tím, že~sám~už~pracovat nemůže. \\
\textschl{Sestra Markéta --} o bratra se jako jediná stará, opravdu se snaží.
\textithl{Cítí to jako povinnost} se o něj starat i v jeho stavu.
Ale čas i~jí~otupí a~i~jí~se~stává přítěží. \\
\textschl{Služka --} najatá služka, o Řehoře se stará,
kdy už nikdo jiný do jeho pokoje nechodí.
Stará se o něj až roboticky, ale stará.
Je to i ona, kdo \textithl{odklidí jeho mrtvé tělo}. \\
\textschl{Prokurista --} pracovník z Řehořovi staré práce,
který se za ním hned první den staví, aby ho přinutil jít pracovat.
Od svého cíle ustupuje, když~se~přesvědčí o~jeho situaci. \\
\textschl{Nájemníci --} když Řehořova rodina trpí na nedostatek peněz,
pronajme část svého obydlí třem mužům.
Řehoř je vyděsí a oni utečou bez placení.
Je to i tato událost, která přesvědčí Řehořovu rodinu,
že \textithl{není ničím víc, než přítěží}, které se musí zbavit.

\subwection{Děj:}
\noindent Příběh začíná krátce po Řehořově metamorfóze v obludného škůdce
\texthl{(německy \enquote{ungeheueres Ungeziefer})}.
Řehoř se nestará o~to, co~se~stalo s~jeho tělem.
Proč má \textit{tolik končetin}.
Proč má tak \textit{zvláštní proporce}.
Proč jeho \textit{záda připomínají krunýř}.
Stará se pouze o to, jak se v jeho stavu \textbfhl{dostavit do práce}.
Během následujících pár desítek minut ztrácí i poslední zbytky lidskosti,
kdy už déle ani \textbfhl{není schopen komunikovat}.
Jeho rodina a prokurista ho nalézají v jeho momentálním stavu.
Během následujících dní se otec musí vrátit do své práce,
sestra se stará o Řehoře a matka lamentuje nad neštěstím, co je postihlo.
Není schopna se podívat na to, co se z jejího syna stalo.
To vede i k situaci, že když se Řehoř jednou omylem nechá zahlédnout,
tak matka omdlí a otec, myslíc si, že se Řehoř dopustil na matce násilí,
\texthl{syna ošklivě poraní}.
Aby vydělali peníze, pronajmou část bytu třem pánům.
Jednoho večera jim Markéta hraje hudbu a Řehoř se musí jít podívat.
Vyděsí nájemníky a ti utečou bez placení, což \texthl{rodinu ještě více obrátí proti Řehořovi}.
Rodině se začíná dařit i bez Řehoře.
Na předchozí hlavu rodiny začínají \textithl{vzpomínat se záští}.
Najmou se dokonce i starou služku a ta se o Řehoře jakž takž stará.
Rána způsobená jablka, jenž po něm otec hodil, se mu nelepší a Řehoř slábne.
Nakonec umírá, služka odklízí jeho tělo a rodina \texthl{spokojeně vykračuje vstříc lepším zítřkům}.

\subwection{Kompozice, prostor a čas:}
\noindent Novela je rozdělena do tří celků.
Navazujících chronologicky na sebe.
První pojednává o \texthl{Řehořově proměně} a reakcích rodiny.
Druhá pak ukazuje \texthl{Řehořovo odloučení} a odcizení rodiny od něj.
Třetí pak vyústí do \texthl{Řehořovi smrti}.
Dílo se odehrává v bytě Samsových,
kdy většina příběhu je~uzavřena pouze mezi čtyřmi stěnami Řehořova pokoje,
což ještě více podtrhuje odloučení a uzavření mladého Samsy.

\subwection{Význam sdělení:}
\noindent Kafkova \textbfhl{Metamorfóza} je jedním z těch příběhů,
co diváka vezme na absurdní cestu do autorova nitra. \\
Jedná se snad o příběh \texthl{muže v depresích}, který se zprotivil sám sobě
a rodina ho pouze využívala?
Či se mladý muž odcizil sám sobě natolik, že~se~přeměnil ve~škůdce
a odcizení rodiny si zavinil sám svým přístupem k životu?
Měla se Řehořova rodina více snažit porozumět jejich synu,
když ani nevěděli, zda to je Řehoř?
Nebo měl Řehoř odejít pryč, do neznáma, a nechat jeho rodinu jejich spokojenému žití?

\wection{LITERÁRNÍ HISTORIE}

\subwection{Politická situace:}
\noindent Rakousko-Uhersko jako státní celek,
složen z patnácti výrazných a mnoho dalších minoritních národností,
čelil značnému \texthl{vnitřnímu napětí} a~\texthl{národnostním rozdílům}.
Ty byly znatelné hlavně v multikulturních místech, kterými byli významná města, kam se řadila i Praha.
20. století zároveň bylo stoletím \texthl{technologického a sociálního pokroku}.
Což lidem vzalo jistoty a vzbuzovalo skepsi.

\subwection{Základní principy fungování společnosti:}
\noindent Praha jako jedno z významnějších měst Rakouska-Uherska, se stalo centrem kultury,
kde se prolínaly vlivy hned z několika směrů.
Kafka byl součástí pražských spisovatelů, mezi kterými převládala pesimismus vůči společnosti.

% \subwection{Kontext dalších druhů umění:}

% \subwection{Kontext literárního vývoje:}

\subwection{Franz Kafka {\ssmall -- život autora:}}
\noindent Syn bohatého židovského obchodníka, se kterým ale \textithl{měl složitý vztah}.
Po vystudování práv se stal úředníkem v pojišťovně, práce ho nenaplňovala,
připadala mu \texthl{byrokratická}, ale dodávala mu životní jistotu.
V jeho dílech se odrážela témata odcizení, izolace a právě absurdity byrokracie.

% \subwection{Vlivy na dané dílo:}

\subwection{Další autorova tvorba:}
\noindent 
Mezi nejznámnější díla patří \textschl{Proces}, \textschl{Zámek}
a \textschl{Amerika} (druhým názvem \textschl{Nezvěstný}).

% \subwection{Jak dílo inspirovalo další vývoj literatury:}

% \wection{LITERÁRNÍ KRITIKA}

% \subwection{Dobové vnímání díla a jeho proměny:}

% \subwection{Aktuálnost tématu a zpracování tématu:}


\wection{OSOBNÍ NÁZOR}
\noindent Příběh byl čtivý. Díky autorově averzi k zodpovídání otázek jsem do konce hltal každou větu,
doufajíc, že dostanu odpověď aspoň na jednu ze svých otázek.
To se sice nestalo, ale to příběhu neubírá na kvalitě.
K mému nesouhlasu to funguje. A funguje to až moc dobře.
Nemělo by, ale~funguje. Zároveň mě zaujalo, jak se tam vlastně nikomu nedá nic vytknout.
Člověk se může zastávat Řehoře a vyčítat jeho rodině jejich přístup k~němu.
Ale oni nevěděli, zda to vůbec je on.
Či mu může člověk vyčítat, že byl pro rodinu, o kterou se dříve staral, stal přítěží,
ale~nic~jiného se~ani~dělat nedalo.

\vfill

\noindent\begin{minipage}{\textwidth}
    {\textcolor{\wpagecolor}{\rule{\linewidth}{0.4pt}}
    \footnotesize
    \textbfhl{Expresionismus --} snaží se vyjádřit vnitřní prožitky a emoce člověka deformací reality;
    \textbfhl{Existencialismus --} zdůrazňuje individuální existenci, svobodu a odpovědnost člověka v absurdním světě;
    \textbfhl{\enquote{Ungeheueres Ungeziefer} --} ohromný\wsl obludný\wsl nestvůrný\wsl strašlivý~
    \textbfhl{\leftrightarrow}{ }~hmyz\wsl havěť\wsl škůdce\wsl breberka
    \textbfhl{$\Rightarrow$}~v~dokumentu jsem se rozhodl pro varianty \texthl{hmyzu} a \texthl{škůdce},
    oproti v češtině používanější variantě o \texthl{brouku}.
    }
\end{minipage}

\newpage

\end{document}